\section{Results}
\label{sec:results}

\subsection{Delegation without monitoring}

\begin{proposition}
\label{prop:keep}
Without monitoring, the manager keeps the $\min\{n, \lfloor A \rfloor\}$ most important tasks and delegates the rest.
\end{proposition}

The intuition is simple. Keeping task $i$ rather than delegating it gains $b v_i$, which is increasing in $v_i$. Attention is therefore spent on the most important tasks first. The proof is in Appendix~\ref{app:proofs}.

\subsection{Delegation with monitoring}

\begin{proposition}
\label{prop:nonmonotone}
With monitoring, the number of delegated tasks is non-monotone in $A$: it is decreasing in $A$ for small $A$ and increasing in $A$ for large $A$.
\end{proposition}

When attention is scarce, every unit is spent on keeping the most valuable tasks. When attention is abundant, a unit of monitoring improves the payoff on every delegated task at once, so the manager prefers to delegate more and monitor better. The proof is in Appendix~\ref{app:proofs}.

\subsection{Implications for firm boundaries}

A task that is delegated with high residual misalignment is a candidate for outsourcing. Proposition~\ref{prop:nonmonotone} implies that firms with intermediate managerial attention integrate the most tasks.
